\documentclass[a4paper,11pt]{article}
\usepackage[margin=1in]{geometry}
\usepackage{algorithm}
\usepackage{algpseudocode}
\usepackage{booktabs}
\usepackage{tabularx}
\usepackage{hyperref}
\usepackage{float}
\usepackage{enumitem}

\title{\textbf{Iteration \#05   Technical Iteration} \\ \large Pseudocode, GitHub Setup, and Implementation Planning \\ \bigskip \normalsize Sensor Fusion System for Autonomous Vehicle Simulation}
\author{Jeongheon Han (David), Charles Wan, Anish Machiraju \\ EECE 2140   Professor Nafa}
\date{April 2026}

\begin{document}
\maketitle

\noindent\textbf{GitHub Repository:} \url{https://github.com/anishmachiraju/sensor-fusion-system}

%% ============================================================
\section{Detailed Pseudocode}
%% ============================================================

The pseudocode below covers five major functionalities of the sensor fusion system. Each section is written clearly enough to be translated directly into C++ code.

\begin{table}[H]
\centering
\small
\renewcommand{\arraystretch}{1.35}
\begin{tabularx}{\textwidth}{>{\bfseries}p{3cm} X p{3.2cm} p{3.2cm}}
\toprule
Functionality & Purpose & Input & Output \\
\midrule
Generate Sensor Reading &
Produces a noisy measurement from the true vehicle state for a given sensor type. &
True vehicle state, noise parameters &
Noisy \texttt{StateEstimate} \\
\addlinespace
Fuse Sensor Data &
Combines calibrated readings from all registered sensors into one accurate estimate using weighted averaging. &
Vector of sensor readings, weight vector &
Fused \texttt{StateEstimate} \\
\addlinespace
Step Vehicle Simulation &
Advances the vehicle's true state by one timestep using a kinematic motion model. &
Current state, speed, turn rate, $dt$ &
Updated true state \\
\addlinespace
Predict Next State &
Projects the current fused estimate forward in time for look-ahead planning. &
Current fused state, $dt$ &
Predicted \texttt{StateEstimate} \\
\addlinespace
Main Simulation Loop &
Orchestrates the full pipeline each timestep: step vehicle, generate readings, fuse, update, log, and display. &
Duration, $dt$, initial conditions, sensor config &
Complete state history \\
\bottomrule
\end{tabularx}
\caption{Summary of pseudocode functionalities.}
\end{table}

% ==================== PSEUDOCODE 1 ====================
\begin{algorithm}[H]
\caption{Sensor::generateReading   Generate a Noisy Sensor Reading}
\begin{algorithmic}[1]

\State \textbf{Class} \textit{Sensor} (base class)

\State \hspace{0.5cm} \textbf{Private:}
\State \hspace{1cm} \textit{sensorType} $\rightarrow$ string
\State \hspace{1cm} \textit{noiseSigma} $\rightarrow$ double
\State \hspace{1cm} \textit{biasOffset} $\rightarrow$ double
\State \hspace{1cm} \textit{lastReading} $\rightarrow$ StateEstimate
\State \hspace{1cm} \textit{isCalibrated} $\rightarrow$ bool
\State \hspace{1cm} \textit{calibrationOffset} $\rightarrow$ double
\State \hspace{1cm} \textit{calibrationScale} $\rightarrow$ double

\State \hspace{0.5cm} \textbf{Protected:}
\State \hspace{1cm} \textit{updateRate} $\rightarrow$ double
\State \hspace{1cm} \textit{lastUpdateTime} $\rightarrow$ double

\State \hspace{0.5cm} \textbf{Public:}

\vspace{0.2cm}
\State \hspace{1cm} \textbf{Function} shouldUpdate(currentTime) $\rightarrow$ bool
\State \hspace{1.5cm} interval $\leftarrow$ 1.0 / updateRate
\State \hspace{1.5cm} \If{currentTime $-$ lastUpdateTime $\geq$ interval}
\State \hspace{2cm} \Return true
\State \hspace{1.5cm} \Else
\State \hspace{2cm} \Return false
\State \hspace{1.5cm} \EndIf
\State \hspace{1cm} \textbf{End Function}

\vspace{0.2cm}
\State \hspace{1cm} \textbf{Function} generateReading(trueState) $\rightarrow$ StateEstimate
\State \hspace{1.5cm} \textbf{Validate:} trueState timestamp $\geq 0$
\State \hspace{1.5cm} reading $\leftarrow$ copy of trueState
\State \hspace{1.5cm} noise $\leftarrow$ random Gaussian(mean=0, stddev=noiseSigma)
\State \hspace{1.5cm} Add noise + biasOffset to the relevant fields of reading
\State \hspace{1.5cm} \If{isCalibrated}
\State \hspace{2cm} reading fields $\leftarrow$ (reading fields $-$ calibrationOffset) $\times$ calibrationScale
\State \hspace{1.5cm} \EndIf
\State \hspace{1.5cm} lastReading $\leftarrow$ reading
\State \hspace{1.5cm} lastUpdateTime $\leftarrow$ trueState.timestamp
\State \hspace{1.5cm} \Return reading
\State \hspace{1cm} \textbf{End Function}

\vspace{0.2cm}
\State \hspace{1cm} \textbf{Function} calibrate(offset, scale) $\rightarrow$ void
\State \hspace{1.5cm} \textbf{Validate:} scale $\neq$ 0
\State \hspace{1.5cm} calibrationOffset $\leftarrow$ offset
\State \hspace{1.5cm} calibrationScale $\leftarrow$ scale
\State \hspace{1.5cm} isCalibrated $\leftarrow$ true
\State \hspace{1cm} \textbf{End Function}

\State \textbf{End Class}

\end{algorithmic}
\end{algorithm}

\noindent\textit{Note:} Each sensor subclass (\texttt{IMUSensor}, \texttt{GPSSensor}, \texttt{TemperatureSensor}, \texttt{WheelEncoderSensor}) overrides \texttt{generateReading()} to add noise only to its relevant fields. For example, \texttt{GPSSensor} perturbs only \texttt{positionX} and \texttt{positionY}, and may randomly return a stale reading to simulate signal dropout.

% ==================== PSEUDOCODE 2 ====================
\begin{algorithm}[H]
\caption{GPSSensor::generateReading   Subclass Override with Signal Dropout}
\begin{algorithmic}[1]

\State \textbf{Class} \textit{GPSSensor} \textbf{extends} \textit{Sensor}

\State \hspace{0.5cm} \textbf{Private:}
\State \hspace{1cm} \textit{positionNoiseSigma} $\rightarrow$ double
\State \hspace{1cm} \textit{positionBias} $\rightarrow$ double
\State \hspace{1cm} \textit{dropoutRate} $\rightarrow$ double \Comment{probability of signal loss}

\State \hspace{0.5cm} \textbf{Public:}

\State \hspace{1cm} \textbf{Function} generateReading(trueState) $\rightarrow$ StateEstimate
\State \hspace{1.5cm} \textbf{Validate:} trueState is not empty
\State \hspace{1.5cm} \If{random uniform(0,\,1) $<$ dropoutRate}
\State \hspace{2cm} \Return lastReading \Comment{simulate signal loss}
\State \hspace{1.5cm} \EndIf
\State \hspace{1.5cm} reading $\leftarrow$ copy of trueState
\State \hspace{1.5cm} noiseX $\leftarrow$ Gaussian(0, positionNoiseSigma)
\State \hspace{1.5cm} noiseY $\leftarrow$ Gaussian(0, positionNoiseSigma)
\State \hspace{1.5cm} reading.positionX $\leftarrow$ reading.positionX + noiseX + positionBias
\State \hspace{1.5cm} reading.positionY $\leftarrow$ reading.positionY + noiseY + positionBias
\State \hspace{1.5cm} \If{isCalibrated}
\State \hspace{2cm} Apply calibration offset and scale to position fields
\State \hspace{1.5cm} \EndIf
\State \hspace{1.5cm} lastReading $\leftarrow$ reading
\State \hspace{1.5cm} lastUpdateTime $\leftarrow$ trueState.timestamp
\State \hspace{1.5cm} \Return reading
\State \hspace{1cm} \textbf{End Function}

\State \textbf{End Class}

\end{algorithmic}
\end{algorithm}

% ==================== PSEUDOCODE 3 ====================
\begin{algorithm}[H]
\caption{FusionSystem::fuseData   Weighted-Average Sensor Fusion}
\begin{algorithmic}[1]

\State \textbf{Class} \textit{FusionSystem}

\State \hspace{0.5cm} \textbf{Private:}
\State \hspace{1cm} \textit{sensors} $\rightarrow$ vector of Sensor pointers
\State \hspace{1cm} \textit{weights} $\rightarrow$ vector of doubles
\State \hspace{1cm} \textit{fusedState} $\rightarrow$ StateEstimate

\State \hspace{0.5cm} \textbf{Public:}

\State \hspace{1cm} \textbf{Function} addSensor(sensor, weight) $\rightarrow$ void
\State \hspace{1.5cm} \textbf{Validate:} sensor is not null, weight $> 0$
\State \hspace{1.5cm} Append sensor to sensors list
\State \hspace{1.5cm} Append weight to weights list
\State \hspace{1cm} \textbf{End Function}

\vspace{0.2cm}
\State \hspace{1cm} \textbf{Function} fuseData() $\rightarrow$ StateEstimate
\State \hspace{1.5cm} \textbf{Validate:} sensors list is not empty
\State \hspace{1.5cm} totalWeight $\leftarrow 0$
\State \hspace{1.5cm} fused $\leftarrow$ new StateEstimate with all fields set to 0

\vspace{0.15cm}
\State \hspace{1.5cm} \For{$i \leftarrow 0$ \textbf{to} sensors.size()$-1$}
\State \hspace{2cm} reading $\leftarrow$ sensors[$i$].getLastReading()
\State \hspace{2cm} w $\leftarrow$ weights[$i$]
\State \hspace{2cm} fused.positionX $\leftarrow$ fused.positionX + w $\times$ reading.positionX
\State \hspace{2cm} fused.positionY $\leftarrow$ fused.positionY + w $\times$ reading.positionY
\State \hspace{2cm} fused.velocityX $\leftarrow$ fused.velocityX + w $\times$ reading.velocityX
\State \hspace{2cm} fused.velocityY $\leftarrow$ fused.velocityY + w $\times$ reading.velocityY
\State \hspace{2cm} fused.heading $\leftarrow$ fused.heading + w $\times$ reading.heading
\State \hspace{2cm} fused.temperature $\leftarrow$ fused.temperature + w $\times$ reading.temperature
\State \hspace{2cm} totalWeight $\leftarrow$ totalWeight + w
\State \hspace{1.5cm} \EndFor

\vspace{0.15cm}
\State \hspace{1.5cm} \textbf{Validate:} totalWeight $> 0$
\State \hspace{1.5cm} Divide every fused field by totalWeight \Comment{normalize}
\State \hspace{1.5cm} fusedState $\leftarrow$ fused
\State \hspace{1.5cm} \Return fusedState
\State \hspace{1cm} \textbf{End Function}

\State \textbf{End Class}

\end{algorithmic}
\end{algorithm}

% ==================== PSEUDOCODE 4 ====================
\begin{algorithm}[H]
\caption{Vehicle   Step Simulation, Predict, Update, and Display}
\begin{algorithmic}[1]

\State \textbf{Class} \textit{Vehicle}

\State \hspace{0.5cm} \textbf{Private:}
\State \hspace{1cm} \textit{state} $\rightarrow$ StateEstimate
\State \hspace{1cm} \textit{stateHistory} $\rightarrow$ vector of StateEstimate
\State \hspace{1cm} \textit{speed} $\rightarrow$ double
\State \hspace{1cm} \textit{turnRate} $\rightarrow$ double

\State \hspace{0.5cm} \textbf{Public:}

\State \hspace{1cm} \textbf{Function} stepSimulation(dt) $\rightarrow$ void
\State \hspace{1.5cm} \textbf{Validate:} dt $> 0$
\State \hspace{1.5cm} state.heading $\leftarrow$ state.heading + turnRate $\times$ dt
\State \hspace{1.5cm} \textbf{Wrap} state.heading into range [0, 360)
\State \hspace{1.5cm} headingRad $\leftarrow$ state.heading $\times$ $\pi$ / 180
\State \hspace{1.5cm} state.velocityX $\leftarrow$ speed $\times$ cos(headingRad)
\State \hspace{1.5cm} state.velocityY $\leftarrow$ speed $\times$ sin(headingRad)
\State \hspace{1.5cm} state.positionX $\leftarrow$ state.positionX + state.velocityX $\times$ dt
\State \hspace{1.5cm} state.positionY $\leftarrow$ state.positionY + state.velocityY $\times$ dt
\State \hspace{1.5cm} state.timestamp $\leftarrow$ state.timestamp + dt
\State \hspace{1.5cm} Append state to stateHistory
\State \hspace{1cm} \textbf{End Function}

\vspace{0.3cm}

\State \hspace{1cm} \textbf{Function} predict(dt) $\rightarrow$ StateEstimate
\State \hspace{1.5cm} \textbf{Validate:} dt $> 0$
\State \hspace{1.5cm} predicted $\leftarrow$ copy of state
\State \hspace{1.5cm} predicted.heading $\leftarrow$ predicted.heading + turnRate $\times$ dt
\State \hspace{1.5cm} headingRad $\leftarrow$ predicted.heading $\times$ $\pi$ / 180
\State \hspace{1.5cm} predicted.velocityX $\leftarrow$ speed $\times$ cos(headingRad)
\State \hspace{1.5cm} predicted.velocityY $\leftarrow$ speed $\times$ sin(headingRad)
\State \hspace{1.5cm} predicted.positionX $\leftarrow$ predicted.positionX + predicted.velocityX $\times$ dt
\State \hspace{1.5cm} predicted.positionY $\leftarrow$ predicted.positionY + predicted.velocityY $\times$ dt
\State \hspace{1.5cm} predicted.timestamp $\leftarrow$ predicted.timestamp + dt
\State \hspace{1.5cm} \Return predicted
\State \hspace{1cm} \textbf{End Function}

\vspace{0.3cm}

\State \hspace{1cm} \textbf{Function} updateState(fusedEstimate) $\rightarrow$ void
\State \hspace{1.5cm} state $\leftarrow$ fusedEstimate
\State \hspace{1cm} \textbf{End Function}

\vspace{0.3cm}

\State \hspace{1cm} \textbf{Function} displayState() $\rightarrow$ void
\State \hspace{1.5cm} Output ``Time:'', state.timestamp
\State \hspace{1.5cm} Output ``Position: ('', state.positionX, ``,'', state.positionY, ``)''
\State \hspace{1.5cm} Output ``Velocity: ('', state.velocityX, ``,'', state.velocityY, ``)''
\State \hspace{1.5cm} Output ``Heading:'', state.heading, ``degrees''
\State \hspace{1.5cm} Output ``Temperature:'', state.temperature, ``C''
\State \hspace{1cm} \textbf{End Function}

\State \textbf{End Class}

\end{algorithmic}
\end{algorithm}

% ==================== PSEUDOCODE 5 ====================
\begin{algorithm}[H]
\caption{Main Simulation Loop   Full Pipeline}
\begin{algorithmic}[1]

\State \textbf{Function} runSimulation(duration, dt, initialState)

\vspace{0.15cm}
\State \Comment{  Validate configuration  }
\State \textbf{Validate:} duration $> 0$, dt $> 0$, dt $\leq$ duration

\vspace{0.15cm}
\State \Comment{  Create vehicle  }
\State vehicle $\leftarrow$ new Vehicle(initialState)

\vspace{0.15cm}
\State \Comment{  Create and calibrate sensors  }
\State imu $\leftarrow$ new IMUSensor(accelSigma, gyroBias, imuRate)
\State gps $\leftarrow$ new GPSSensor(posSigma, posBias, gpsRate)
\State temp $\leftarrow$ new TemperatureSensor(tempSigma, tempBias, tempRate)
\State encoder $\leftarrow$ new WheelEncoderSensor(tickSigma, tickBias, radius, ticks)
\State imu.calibrate(imuOffset, imuScale)
\State gps.calibrate(gpsOffset, gpsScale)
\State temp.calibrate(tempOffset, tempScale)
\State encoder.calibrate(encOffset, encScale)

\vspace{0.15cm}
\State \Comment{  Register sensors with fusion system  }
\State fusion $\leftarrow$ new FusionSystem()
\State fusion.addSensor(imu, 0.30)
\State fusion.addSensor(gps, 0.35)
\State fusion.addSensor(temp, 0.10)
\State fusion.addSensor(encoder, 0.25)

\vspace{0.15cm}
\State \Comment{  Run simulation  }
\State currentTime $\leftarrow 0$
\While{currentTime $<$ duration}
\State \hspace{0.5cm} vehicle.stepSimulation(dt) \Comment{advance ground truth}
\State \hspace{0.5cm} trueState $\leftarrow$ vehicle.getState()
\State \hspace{0.5cm} \For{each sensor in [imu, gps, temp, encoder]}
\State \hspace{1cm} \If{sensor.shouldUpdate(currentTime)}
\State \hspace{1.5cm} sensor.generateReading(trueState)
\State \hspace{1cm} \EndIf
\State \hspace{0.5cm} \EndFor
\State \hspace{0.5cm} \textbf{Try:}
\State \hspace{1cm} fusedEstimate $\leftarrow$ fusion.fuseData()
\State \hspace{1cm} vehicle.updateState(fusedEstimate)
\State \hspace{1cm} vehicle.displayState()
\State \hspace{0.5cm} \textbf{Catch} exception:
\State \hspace{1cm} Output error message, continue to next step
\State \hspace{0.5cm} currentTime $\leftarrow$ currentTime + dt
\EndWhile

\vspace{0.15cm}
\State \Comment{  Final output  }
\State history $\leftarrow$ vehicle.getHistory()
\State Output formatted table of all states in history
\State \Return history

\State \textbf{End Function}

\end{algorithmic}
\end{algorithm}

%% ============================================================
\section{Class Interaction and Responsibility}
%% ============================================================

\begin{table}[H]
\centering
\small
\renewcommand{\arraystretch}{1.4}
\begin{tabularx}{\textwidth}{>{\bfseries\ttfamily}p{2.2cm} X p{2.8cm} X}
\toprule
\textnormal{\textbf{Class Name}} & Responsibility & Interacts With & Reason \\
\midrule
StateEstimate &
Stores all vehicle state variables (position, velocity, heading, temperature, timestamp) in a single lightweight struct. &
Sensor, FusionSystem, Vehicle &
Provides a standardized data type so every component can exchange state information without ambiguity. \\
\addlinespace
Sensor \textnormal{(base)} &
Defines the shared interface for all sensor types: stores noise and calibration parameters, generates noisy readings, and tracks update timing. &
StateEstimate, FusionSystem &
Encapsulates common sensor behavior; enables polymorphism so FusionSystem can treat every sensor uniformly through a base pointer. \\
\addlinespace
IMUSensor &
Generates noisy acceleration and heading data by overriding \texttt{generateReading()} to perturb only inertial fields. &
Sensor (inherits), StateEstimate &
Models the inertial measurement unit, critical for estimating orientation and short-term velocity changes. \\
\addlinespace
GPSSensor &
Generates noisy position data and simulates occasional signal dropout via a configurable dropout rate. &
Sensor (inherits), StateEstimate &
Models GPS, providing absolute position but subject to latency and satellite occlusion. \\
\addlinespace
Temperature- Sensor &
Generates noisy temperature readings with a simulated thermal response delay. &
Sensor (inherits), StateEstimate &
Models an ambient temperature sensor for environmental monitoring. \\
\addlinespace
WheelEncoder- Sensor &
Generates noisy velocity estimates based on wheel rotation using wheel radius and ticks per revolution. &
Sensor (inherits), StateEstimate &
Models odometry input, providing continuous velocity data that complements GPS and IMU. \\
\addlinespace
FusionSystem &
Registers multiple sensors, collects their latest calibrated readings, and computes a weighted-average fused state estimate. &
Sensor (via pointers), StateEstimate &
Central intelligence of the system; no single sensor is reliable alone, so fusion combines them for accuracy. \\
\addlinespace
Vehicle &
Maintains the ground-truth state, advances it using kinematics, stores the full state history, and displays output to console. &
StateEstimate, FusionSystem (receives fused estimate) &
Represents the autonomous vehicle being simulated; owns the truth that sensors observe and the log that captures results. \\
\bottomrule
\end{tabularx}
\caption{Class interaction and responsibility overview.}
\end{table}

\subsection{How the Classes Work Together}

The simulation follows a clear chain of responsibility:

\begin{enumerate}[leftmargin=1.5cm]
\item \textbf{Vehicle} holds the ground-truth state and advances it each timestep using a kinematic model (\texttt{stepSimulation}).
\item Each \textbf{Sensor} subclass (\texttt{IMUSensor}, \texttt{GPSSensor}, \texttt{TemperatureSensor}, \texttt{WheelEncoderSensor}) observes the true state through its \texttt{generateReading()} method, adding type-specific noise and bias. Polymorphism allows the \texttt{FusionSystem} to call the same interface on any sensor.
\item \textbf{FusionSystem} iterates over all registered sensors, retrieves their latest readings via \texttt{getLastReading()}, and computes a weighted-average \texttt{StateEstimate}. The fused result is more accurate than any individual sensor.
\item The fused estimate is fed back to \textbf{Vehicle} via \texttt{updateState()}, simulating how a real autonomous vehicle would correct its self-estimate based on sensor fusion.
\item \textbf{StateEstimate} is the shared data contract used in every exchange, ensuring type safety and clarity across all classes.
\end{enumerate}

%% ============================================================
\section{GitHub Repository Structure}
%% ============================================================

The repository has been updated using Git command-line tools (\texttt{git add .}, \texttt{git commit -m "..."}, \texttt{git push}). The current structure is:

\begin{verbatim}
sensor-fusion-system/
|-- README.md
|-- docs/
|   |-- System_Design_Overview.pdf
|   |-- Iteration05_Pseudocode.pdf
|-- pseudocode/
|   |-- pseudocode.txt
|-- src/
|   |-- main.cpp
|   |-- StateEstimate.h
|   |-- Sensor.h
|   |-- Sensor.cpp
|   |-- IMUSensor.h
|   |-- IMUSensor.cpp
|   |-- GPSSensor.h
|   |-- GPSSensor.cpp
|   |-- TemperatureSensor.h
|   |-- TemperatureSensor.cpp
|   |-- WheelEncoderSensor.h
|   |-- WheelEncoderSensor.cpp
|   |-- FusionSystem.h
|   |-- FusionSystem.cpp
|   |-- Vehicle.h
|   |-- Vehicle.cpp
|-- images/
    |-- system_diagram.png
\end{verbatim}

%% ============================================================
\section{README.md Summary}
%% ============================================================

The repository README.md includes:

\begin{itemize}[leftmargin=1.5cm]
\item \textbf{Project Title:} Sensor Fusion System for Autonomous Vehicle Simulation
\item \textbf{Team Members:} Jeongheon Han (David), Charles Wan, Anish Machiraju
\item \textbf{Project Overview:} A C++ simulation that integrates data from IMU, GPS, temperature, and wheel encoder sensors using a weighted-average fusion algorithm to estimate and predict the motion and state of an autonomous vehicle.
\item \textbf{Main Functionalities:} Sensor simulation with noise, calibration, weighted-average data fusion, kinematic state prediction, simulation loop, formatted output, and exception handling.
\item \textbf{OOP Design Summary:} Eight components \texttt{StateEstimate} struct, \texttt{Sensor} base class, four sensor subclasses, \texttt{FusionSystem}, and \texttt{Vehicle} connected via polymorphism and a shared data contract.
\item \textbf{Tools and Technologies:} C++17, Git / GitHub, VS Code, Google Workspace, Claude.
\item \textbf{Folder Structure:} \texttt{src/} for source code, \texttt{docs/} for PDFs, \texttt{pseudocode/} for planning, \texttt{images/} for diagrams.
\item \textbf{Build Instructions:} \texttt{g++ -std=c++17 -o simulation src/*.cpp} then \texttt{./simulation}.
\item \textbf{Project Goals:} Demonstrate sensor fusion concepts, practice collaborative C++ development with Git, and produce a well-documented, fully functional simulation.
\end{itemize}

\end{document}
